% !TeX root = Chapter_04_Risk_Quantification_Qualification.tex
% Chapter 4 — Risk Quantification and Qualification
\documentclass[11pt,openany]{book}
\usepackage{graphicx}
\usepackage{amsmath,amssymb}
\usepackage{hanging}
\makeatletter
\def\input@path{{second_ed/style/}{../style/}}
\makeatother
\usepackage{erm_textbook_style}
\providecommand{\tightlist}{\setlength{\itemsep}{0pt}\setlength{\parskip}{0pt}}
\emergencystretch=2em

% Simple display for Excel formulas and short code snippets.
% Use \excelcode{formula} for inline formulas and the excelblock
% environment for multi-line blocks (see style guide Section 18).
\newcommand{\excelcode}[1]{{\small\ttfamily #1}}
\newenvironment{excelblock}%
  {\begin{quote}\small\ttfamily\setlength{\parskip}{2pt}}%
  {\end{quote}}

\begin{document}

\setcounter{chapter}{3}
\chapter{Risk Quantification and Qualification}
\label{chap:risk-quantification}

% -----------------------------------------------------------------------
\begin{learningobjectives}
  \item Distinguish between qualitative and quantitative approaches to risk
        assessment and explain when each is most appropriate.
  \item Explain why firms cannot rely solely on intuition to understand their
        major risks and the value of systematic measurement.
  \item Compute and interpret basic risk metrics including frequency, severity,
        expected loss, variance, and standard deviation from sample data.
  \item Describe how a retailer and a manufacturer might systematically
        quantify their key operational and hazard risks.
  \item Recognize and interpret normal versus skewed loss distributions,
        including the concept of fat tails and tail risk.
  \item Use Excel tools to estimate the probability and severity distribution
        of losses and calculate a Value at Risk (VaR) measure.
  \item Explain why VaR is useful but incomplete as a risk measure and why
        qualitative judgment remains essential.
  \item Integrate quantitative risk metrics with qualitative assessments to
        support enterprise risk management decisions.
\end{learningobjectives}

% -----------------------------------------------------------------------
\chapteroverview{%
In Chapter~3, we learned how to identify and categorize risks using the
four-quadrant framework and various risk identification techniques. Knowing
that risks exist is essential, but it is only the beginning. Organizations
must also assess the \emph{magnitude} of identified risks---how likely they
are to occur and how severe their consequences might be. This chapter
addresses risk assessment through both quantitative and qualitative
approaches.

Risk quantification uses numbers, statistical methods, and models to measure
risk exposure. By analyzing historical data, organizations can estimate the
frequency of adverse events, the distribution of loss severities, and
aggregate risk exposure. Quantitative measures provide objectivity, enable
comparison across different risks, support capital allocation decisions, and
satisfy regulatory and stakeholder expectations for rigorous risk assessment.
However, quantification has limitations: it relies on historical data that may
not predict the future, it cannot easily capture rare catastrophic events, and
it requires judgment about which models and assumptions to use.

Risk qualification uses narrative descriptions, expert judgment, scenario
analysis, and visual tools like risk matrices to assess risks. Qualitative
approaches are valuable when historical data is limited, when risks are complex
and difficult to model, when stakeholder perceptions matter, and when the
organization needs to communicate about risks to non-technical audiences. This
chapter develops skills in both approaches. We begin with the fundamental
statistical concepts underlying risk measurement, work through two detailed
cases---a regional retailer and a parts manufacturer---examine loss
distributions with particular attention to skewed distributions and tail risk,
and build a workers' compensation model in Excel to calculate Value at Risk.
The chapter concludes by addressing the limits of quantitative methods and
showing how to integrate quantitative and qualitative assessments into
enterprise risk management.%
}

\clearpage

% -----------------------------------------------------------------------
% Opening Case 1: Long-Term Capital Management
% -----------------------------------------------------------------------
\begin{corporateexample}[When the Model Was the Risk: Long-Term Capital Management]

In the mid-1990s, Long-Term Capital Management (LTCM)\index{Long-Term Capital Management (LTCM)} was the most admired
hedge fund on Wall Street. Founded in 1994 by John Meriwether\index{Meriwether, John}, formerly of
Salomon Brothers, the firm assembled a team that included two Nobel
Prize-winning economists---Myron Scholes\index{Scholes, Myron} and Robert Merton\index{Merton, Robert}---along with some
of the best fixed-income traders and quantitative analysts in the world.
Investors lined up to get in. The fund charged high fees, limited withdrawals,
and cultivated an aura of exclusivity.

LTCM's strategy was ``relative-value arbitrage.'' Rather than betting on the
overall direction of markets, it looked for tiny mispricings between similar
securities---say, two government bonds with nearly identical cash flows that,
in theory, should trade at nearly the same yield. The fund would buy the
cheaper bond and short the more expensive one, expecting the spread to
converge. The individual trades had small expected profits, but LTCM multiplied
those profits through enormous leverage\index{leverage}: it borrowed heavily and used
derivatives so that a dollar of investor capital controlled many dollars of
positions.

Quantification was central to the strategy. LTCM's models used historical data
to estimate how volatile different spreads were, how correlated they were with
one another, and how likely it was that many positions would move against the
fund simultaneously. Using techniques that would later be popularized as value
at risk (VaR)\index{Value at Risk (VaR)}, the firm calculated measures like ``on 99 out of 100 days, our
losses should not exceed \$X'' or ``the probability of losing more than \$Y in
a month is less than 1 in 10,000.'' These numbers reassured investors and
justified the high leverage. If the worst-case loss was extremely unlikely, why
not borrow more to scale up the small arbitrage profits?

For a few years, the models seemed to work. Between 1994 and 1997, LTCM
generated very high returns with remarkably low reported volatility. Its VaR
numbers suggested the firm was taking far less risk than most funds holding
equities. Regulators and counterparties took comfort from the sophisticated
math and the distinguished people running it. On paper, this looked like a
textbook case of quantification done right.

In 1998, the real world refused to cooperate.

That summer, Russia defaulted on its government debt and devalued the ruble.
Investors panicked. Instead of behaving independently, prices across many
different securities moved together in a ``flight to quality'': riskier assets
were sold indiscriminately and the safest government bonds attracted all the
demand. The spreads LTCM had bet would converge suddenly blew out. Positions
that were supposed to offset one another all lost value at the same time.
Liquidity dried up; the fund could not exit its positions without pushing
prices even further against itself.

The losses escalated quickly. Over a few months in 1998, LTCM lost about
\$4.6~billion---roughly 90~percent of its capital. Trades that the models said
should almost never produce large losses did exactly that. Correlations that
had been low or negative in the historical data spiked toward one. The
``six-sigma'' or ``ten-sigma'' scenarios treated as virtually impossible turned
out to be quite possible in a crisis.

The problem was not that LTCM failed to quantify risk. It quantified risk
extensively. The failure lay in what exactly was being quantified and what
assumptions were embedded in the numbers. The VaR models assumed that the past
was a reasonable guide to the future, that return distributions were fairly
stable, and that diversification across many trades would protect the fund
because not everything would go wrong at once. They also largely ignored
liquidity risk\index{liquidity risk}---the possibility that in a stressed market, LTCM would be
unable to exit positions without enormous price impact. When those assumptions
broke down, the reassuring numbers became dangerous. High leverage magnified
the error: a position sized to be ``safe'' given a 99~percent VaR can be
catastrophic if the true distribution has fatter tails than the model allows.
LTCM's near-failure was large enough, and its links to major banks so
extensive, that the Federal Reserve Bank of New York\index{Federal Reserve Bank of New York} convened a group of large
financial institutions to organize a private rescue and avoid broader market
disruption.

For students of risk quantification, LTCM offers several lessons. Sophisticated
models do not eliminate risk; they transform model risk\index{model risk} into a central risk
factor. Quantifying risk with historical data is most fragile in exactly the
tail events that matter most. Leverage interacts dangerously with overconfident
quantification: if you size positions based on optimistic estimates of
volatility and correlation, you are effectively betting your firm on the model
being right. The LTCM episode illustrates why enterprise risk management\index{enterprise risk management (ERM)} must
ask not only ``What does the VaR say?'' but also ``What assumptions does this
number rest on, and how might they fail under stress?''

\source{Lowenstein, R. (2000). \emph{When genius failed: The rise and fall of
Long-Term Capital Management.} Random House.}

\end{corporateexample}

% -----------------------------------------------------------------------
% Opening Case 2: Moneyball
% -----------------------------------------------------------------------
\begin{corporateexample}[Moneyball: Quantification as Competitive Advantage]

\begin{quote}
\textit{``This is a very simple game. You throw the ball, you catch the ball,
you hit the ball. Sometimes you win, sometimes you lose, sometimes it rains.''}\\[2pt]
\hfill ---Ebby Calvin ``Nuke'' LaLoosh, \textit{Bull Durham} (1988)
\end{quote}

\medskip
Wouldn't it be nice to reduce some of that uncertainty?

At the start of the 2002 Major League Baseball season, the Oakland Athletics\index{Oakland Athletics}
had a problem. They had just lost three star players---Jason Giambi, Johnny
Damon, and Jason Isringhausen---to richer teams in free agency. The New York
Yankees, Boston Red Sox, and other big-market clubs could afford to pay top
dollar for proven talent. The A's could not. Their payroll was less than a
third of the Yankees', and there was no salary cap to level the playing field.

Traditional baseball wisdom said that without stars, a small-market team had
little chance. Scouts and executives judged players with a mix of experience,
intuition, and conventional statistics like batting average, runs batted in
(RBIs), and pitcher wins. These measures were familiar, but they were also
noisy and sometimes misleading. A player might have a high batting average
because of a hot streak or lucky singles; a pitcher might accumulate wins
because his team scored a lot of runs, not because he actually prevented runs
better than others.

Oakland's general manager, Billy Beane\index{Beane, Billy}, and his assistant Paul DePodesta\index{DePodesta, Paul} took
a different approach. Influenced by sabermetric research\index{sabermetrics}---most famously the
work of Bill James\index{James, Bill}---they set out to quantify what actually led to wins, and to
measure players accordingly. They treated the problem like an applied
risk-quantification exercise: given limited resources, how could they buy as
much expected run production and run prevention as possible, while managing the
risk that their measurements were wrong?

Their analysis reached several conclusions that contradicted conventional
wisdom. One of the most important was that on-base percentage
(OBP)\index{on-base percentage (OBP)}---how often a player reached base by any means, including
walks---was far more valuable than teams were paying for. Walks, which scouts
often dismissed as signs of passivity, turned out to correlate strongly with
run scoring. Another finding was that certain types of players---older hitters
with high OBP but unconventional bodies or ``ugly'' swings---were
systematically undervalued in the market.

In Moneyball\index{Moneyball} terms, the A's quantified ``player risk'' differently. Instead of
asking ``Does this player look like a star?'' they built simple models
estimating how many runs and wins a player was likely to add, given his
underlying statistics. They also considered variance: a slugger with a lot of
strikeouts might have more volatile performance than a high-OBP player, even if
both were projected to contribute similar runs. For a low-budget team, that
volatility was itself a risk---there was little room to absorb a prolonged
slump from an expensive acquisition.

Armed with these quantifications, the A's made unconventional choices. They
signed players like Scott Hatteberg, a catcher with arm problems whom they
converted to first base, largely because he reached base at a high rate but
other teams viewed him as damaged goods. They avoided paying for traditional
statistics the market overvalued, such as RBIs, which depend heavily on context
rather than individual skill.

The results were striking. In 2002, despite their low payroll, the A's won 103
games---matching the Yankees---and reached the playoffs. They did it not by
outspending rivals, but by re-pricing risk. Players the market treated as risky
or marginal were, by the A's calculations, bargains with relatively predictable
on-base skills. Conversely, many high-priced free agents were overpriced
lottery tickets.

The A's models were not perfect. Some bets failed. Injuries and random
variation still mattered. Over time, other teams copied the approach, eroding
Oakland's advantage. But the Moneyball story illustrates something that applies
well beyond baseball: better quantification of risk and value changes which
opportunities you see and which you ignore. If you measure only traditional
metrics, you may systematically overpay for familiar ``stars'' and underinvest
in boring but reliably productive alternatives. Quantification is most powerful
when paired with constraints---the A's did not try to eliminate risk but tried
to choose the mix of players that maximized expected wins per dollar under a
strict budget, much like a firm optimizing a portfolio of investments under
capital limits.

\source{Lewis, M. (2003). \emph{Moneyball: The art of winning an unfair game.}
W.\,W.\ Norton \& Company.}

\end{corporateexample}

% -----------------------------------------------------------------------
\section{Why Quantify and Qualify Risk?}
\label{sec:why-quantify}
% -----------------------------------------------------------------------

Risk assessment\index{risk assessment} is the bridge between risk identification (Chapter~3) and risk
response (future chapters). Once an organization knows what risks it faces, it
must evaluate how significant those risks are. Which risks deserve the most
management attention? Which require substantial resources for mitigation or
transfer? Which fall within the organization's risk appetite, and which exceed
acceptable levels? These questions demand systematic assessment, not mere
intuition.

\subsection{The Limitations of Intuition and Informal Assessment}
\label{subsec:intuition-limits}

Many managers have strong intuitions about risks based on experience, industry
knowledge, and pattern recognition. These intuitions are valuable---experienced
managers often detect risks that purely statistical approaches miss. But
relying solely on intuition creates several problems.

First, \textbf{human judgment is subject to systematic biases.} Psychological
research shows that people tend to overestimate the likelihood of vivid, recent,
or emotionally salient events while underestimating risks that are abstract or
unfamiliar. A manager who recently experienced a cybersecurity incident may
overestimate cyber risk while underestimating equally significant but less
visible risks. People also exhibit overconfidence\index{overconfidence}, believing their judgments are
more accurate than they actually are, and anchoring bias\index{anchoring bias}, being influenced by
initial estimates even when subsequent information suggests different
conclusions.

Second, \textbf{informal assessment makes comparison and prioritization
difficult.} If one manager says a particular risk is ``high'' while another
says a different risk is ``moderate,'' how should the organization allocate
resources between them? Without common metrics and systematic measurement,
resources may flow to whichever risks are advocated most forcefully by
influential managers rather than to risks that are objectively most
significant.

Third, \textbf{stakeholders increasingly expect rigorous, data-driven risk
assessment.} Boards of directors, regulators, rating agencies, investors, and
lenders want evidence that management understands its risk exposures
quantitatively. Regulatory frameworks for banks and insurers mandate
quantitative risk measurement and minimum capital based on risk models. Rating
agencies assess the sophistication of risk quantification when assigning credit
ratings. Informal, intuitive risk assessment no longer satisfies these
expectations.

\subsection{The Value of Quantitative Risk Measurement}
\label{subsec:quantitative-value}

Quantitative risk measurement\index{risk quantification} addresses these limitations by bringing
discipline, objectivity, and comparability to risk assessment. Several specific
benefits warrant the effort.

\textbf{Capital allocation and risk appetite.} Organizations have finite capital
and must decide how much to allocate to different business units, products, or
activities partly based on their riskiness. Risk-adjusted performance
measurement---comparing returns earned relative to risks taken---requires
quantifying risk. A business unit generating 15~percent return on capital might
appear attractive until risk measurement reveals it also involves high
volatility or potential for large losses. Defining risk appetite\index{risk appetite} also requires
quantification: statements like ``We will not accept risks that could reduce
annual earnings by more than 10~percent'' or ``We will maintain capital
sufficient to survive a 1-in-100-year loss event'' require measuring what those
thresholds mean numerically.

\textbf{Pricing and reserving.} Insurance companies must price policies based on
expected losses plus expenses and profit margin. Without quantifying expected
claim frequency and severity, insurers cannot price rationally. Similarly,
manufacturers must set warranty reserves based on expected warranty claim costs,
and retailers must establish loss prevention budgets based on expected
shrinkage. Any business decision that involves absorbing risk requires
estimating the expected cost of that risk.

\textbf{Risk transfer and insurance purchasing.} Organizations deciding how much
insurance to buy and what deductibles or retentions to accept need to understand
their loss distributions quantitatively. An organization might rationally choose
to retain small, frequent losses---because insurance for such losses is
expensive relative to expected losses---while transferring large, infrequent
losses where insurance provides genuine value. Making this decision requires
quantifying the frequency and severity of different loss scenarios.

\textbf{Regulatory compliance.} Banks must calculate risk-weighted assets and
maintain minimum capital ratios under Basel~III\index{Basel III}. These calculations require
quantitative models for credit risk, market risk, and operational risk
(Basel Committee on Banking Supervision, 2011). Insurers face similar requirements under risk-based capital
regulations. Even non-financial firms face regulatory expectations for risk
quantification in areas like environmental liability estimation, pension
liability measurement, and asset impairment testing.

\textbf{Performance monitoring and early warning.} Once risks are quantified,
organizations can establish key risk indicators (KRIs)\index{key risk indicator (KRI)}---metrics that signal
when risk exposure is increasing. If historical analysis shows that workplace
accidents increase when overtime hours exceed certain thresholds, overtime can
serve as a leading indicator of accident risk. Quantitative monitoring enables
proactive risk management rather than reactive response after losses occur.

\subsection{Quantification Beyond Finance: Non-Financial Companies}
\label{subsec:quantification-nonfinancial}

The examples in Section~\ref{subsec:quantitative-value} draw heavily from
banking and insurance, because those industries have faced the most regulatory
pressure to formalize risk measurement. But the underlying logic---measure what
you can, track it consistently, compare against benchmarks---applies just as
well to hospitals, manufacturers, retailers, sports franchises, and virtually
any organization that faces repeated, partially predictable losses.

The Moneyball story illustrates this clearly. Baseball teams had long relied on
intuition and traditional metrics to evaluate players. When the Oakland A's
began asking which statistics actually predicted run scoring---and measured
those statistics rigorously---they found systematic mispricings that the rest
of the market had missed. Their ``risk management'' was simply disciplined
measurement of what mattered, combined with a willingness to act on the numbers
when conventional wisdom pointed elsewhere.

Hospitals offer a less glamorous but equally instructive example. A hospital's
most significant operational risks include hospital-acquired infections (HAIs)\index{hospital-acquired infections},
medication errors, and preventable readmissions---all of which cause patient
harm and generate substantial financial costs. For decades, these were managed
largely through clinical intuition and periodic audits. Beginning in the 2000s,
hospitals began systematically quantifying them. The CDC's National Healthcare
Safety Network (NHSN) tracks central line-associated bloodstream infection
(CLABSI) rates expressed as infections per 1,000 catheter-days, adjusted for
patient acuity, and publishes national benchmarks. A hospital can compare its
rate against the national median, identify which units or procedures drive
excess infections, and set measurable improvement targets. The federal
government made this consequential: under the Affordable Care Act\index{Affordable Care Act}, hospitals
with excess readmission rates for conditions like heart failure and pneumonia
face direct reductions in Medicare reimbursement. A hospital that quantifies
its 30-day readmission rate at 22~percent against a national average of
15~percent for heart failure patients is no longer managing a vague quality
problem---it is managing a measurable financial exposure with identifiable
drivers (inadequate discharge planning, poor patient education, lack of
post-discharge follow-up) that can be addressed with targeted interventions.

The same logic extends to other industries. Airlines track dispatch
reliability---the percentage of flights that depart on schedule without a
mechanical delay or cancellation---as a direct operational risk indicator,
measured per 1,000 departures and compared against industry benchmarks. A
carrier discovering that one aircraft type shows dispatch reliability of
96.8~percent versus a peer average of 98.5~percent knows quantitatively where
it is underperforming and by how much. That number drives maintenance
investment decisions, fleet scheduling, and reserve aircraft positioning in a
way that ``our fleet seems unreliable this year'' simply cannot. What unifies
these examples---baseball, hospitals, airlines---is not analytical
sophistication. None required derivatives pricing or Monte Carlo simulation.
They required identifying a relevant exposure base, consistently measuring
events against that base, tracking severity or impact when events occurred, and
comparing against a benchmark. That discipline, applied consistently, changes
decisions. Without it, organizations tend to respond to the most recent vivid
incident rather than the underlying pattern---exactly the behavioral bias that
systematic quantification is designed to correct.

\subsection{The Continuing Role of Qualitative Assessment}
\label{subsec:qualitative-role}

Despite the value of quantification, qualitative assessment\index{risk qualification} remains essential.
Numbers alone do not provide complete understanding, and some risks resist
quantification entirely.

\textbf{Limited data for low-frequency, high-severity events.} Many significant
risks occur rarely, providing limited historical data. A manufacturer may have
operated for 20~years without experiencing a catastrophic equipment failure, but
that is not evidence the risk is negligible---it may simply mean the rare event
has not yet occurred. Quantifying such risks requires judgment about worst-case
scenarios, not just statistical analysis of historical losses.

\textbf{Emerging and changing risks.} Quantification relies on historical data,
but the future may differ from the past. New technologies create new risks
(cybersecurity risk evolved dramatically as information systems changed).
Regulatory changes create new compliance risks. Strategic risks often involve
unique circumstances without historical precedent. For these risks, qualitative
scenario analysis and expert judgment supplement limited quantitative data.

\textbf{Risk interdependencies and correlations.} Quantifying how different
risks interact---how a natural disaster might simultaneously disrupt supply
chains, damage facilities, and harm employees---requires assumptions about
correlations that are hard to estimate from data. Qualitative analysis through
scenario exercises and stress testing reveals interactions that purely
statistical approaches might miss.

\textbf{Communication and stakeholder engagement.} Most business leaders and
board members are not statisticians. Risk matrices, heat maps\index{heat map}, and narrative
scenario descriptions often communicate more effectively with non-technical
audiences than quantitative models. Moreover, stakeholder perceptions of risk
matter even when they diverge from statistical assessments---a risk that
stakeholders find unacceptable may require attention even if quantitative
analysis suggests it is modest. The goal is informed risk assessment, not
mathematical precision for its own sake.

% -----------------------------------------------------------------------
\section{Basic Statistical Tools for Risk Measurement}
\label{sec:statistical-tools}
% -----------------------------------------------------------------------

Risk quantification builds on fundamental concepts from statistics. This
section reviews the essential tools for measuring risk, assuming you have had
an introductory statistics course but may need a refresher. The emphasis is on
intuitive understanding and practical application, not mathematical derivation.

\subsection{Frequency, Severity, and Expected Loss}
\label{subsec:frequency-severity}

The most basic framework for quantifying risk involves three concepts\index{frequency-severity analysis}:
frequency, severity, and expected loss.

\textbf{Frequency}\index{frequency} refers to how often an adverse event occurs, typically
expressed as the number of events per unit of exposure over a defined time
period. For example:

\begin{itemize}
  \item Workplace accidents per 100 employees per year
  \item Customer slip-and-fall incidents per store per year
  \item Product defects per 10,000 units manufactured
  \item Data breach attempts per month
\end{itemize}

Frequency is estimated from historical data by counting events and dividing by
the exposure base. If a retailer with 50 stores experienced 15 slip-and-fall
incidents over the past year, the estimated frequency is
\(15 / 50 = 0.30\) incidents per store per year.

\textbf{Severity}\index{severity} refers to the magnitude of loss when an event occurs,
measured in dollars, time, or other relevant units. Severity analysis examines
the distribution of loss sizes: how severe is the typical loss? What is the
range from smallest to largest? Are there occasional very large losses that
dominate total losses even though they are infrequent?

\textbf{Expected loss}\index{expected loss} combines frequency and severity to estimate the average
total loss over a defined period:

\[
  \text{Expected Loss} = \text{Frequency} \times \text{Average Severity}
\]

For example, if a retailer expects 0.30 slip-and-fall incidents per store per
year and the average incident costs \$10,000, the expected loss per store per
year is \(0.30 \times \$10{,}000 = \$3{,}000\). For 50 stores, expected annual
loss is \$150,000. Expected loss represents the central tendency---what you
would expect on average over many years. In any single year, actual losses will
differ due to random variation. Expected loss provides the baseline for
budgeting, setting loss reserves, and determining appropriate insurance
coverage.

\subsection{Descriptive Statistics: Mean, Median, Variance, and Standard Deviation}
\label{subsec:descriptive-stats}

Beyond expected loss, several statistical measures help characterize risk
distributions.

\textbf{Mean (average)}\index{mean} is the arithmetic average of observed values. Mean
severity is the average loss size across all events. Means are intuitive and
widely used but can be heavily influenced by outliers---a few very large losses
can pull the mean upward even when most losses are small.

\textbf{Median}\index{median} is the middle value when observations are arranged in order.
Half of observations fall below the median; half above. Median severity is
often considerably lower than mean severity for loss distributions because most
losses are small while a few large losses pull up the mean. If a manufacturer
experiences 100 warranty claims ranging from \$50 to \$50,000, the median might
be \$800 while the mean might be \$2,000, pulled up by a few very large claims.

\textbf{Minimum and maximum} establish the range of historical experience. The
maximum historical loss should not be assumed to be the maximum possible loss:
larger losses than have been experienced historically can and do occur.

\textbf{Variance}\index{variance} measures dispersion around the mean. It is the average squared
deviation from the mean:

\[
  \text{Variance} = \frac{\sum_{i=1}^{n}(x_i - \bar{x})^2}{n}
\]

Variance is expressed in squared units, which is not intuitive, so it is rarely
reported directly. However, it is important conceptually because it quantifies
uncertainty---distributions with higher variance are more uncertain and risky.

\textbf{Standard deviation}\index{standard deviation} is the square root of variance, expressed in the
same units as the original data. For normally distributed data, approximately
68~percent of observations fall within one standard deviation of the mean,
95~percent within two standard deviations, and 99.7~percent within three
standard deviations. If expected annual loss is \$100,000 with standard
deviation \$30,000, actual annual losses will fall between \$70,000 and
\$130,000 in roughly 68~percent of years---assuming approximately normal
distribution. Many risk distributions are not normal, so the standard normal
percentages do not always apply, but standard deviation remains useful as a
general measure of variability.

\subsection{Correlation and Risk Aggregation}
\label{subsec:correlation}

Most organizations face multiple risks simultaneously. Total risk depends not
only on the magnitude of individual risks but also on how those risks relate to
each other---whether they tend to occur together or independently.

\textbf{Correlation}\index{correlation} measures the degree to which two variables move together,
ranging from \(-1\) to \(+1\). Positive correlation (between 0 and \(+1\))
means when one variable increases, the other tends to increase as well.
Workplace accidents and equipment breakdowns might be positively correlated if
both increase during periods of rushed production and inadequate maintenance.
Negative correlation (between 0 and \(-1\)) means when one variable increases,
the other tends to decrease. Heating costs and cooling costs are negatively
correlated---high in winter, low in summer, or vice versa. Zero correlation
means the variables move independently.

Correlation matters for risk aggregation. Consider an organization facing two
risks, each with expected loss of \$1~million and standard deviation of
\$500,000. Total expected loss is always the sum of individual expected losses:
\(\$1\text{M} + \$1\text{M} = \$2\text{M}\), regardless of correlation. Total
standard deviation depends on correlation:

\begin{itemize}
  \item If perfectly positively correlated: \(\$500{,}000 + \$500{,}000 = \$1{,}000{,}000\)
  \item If independent (correlation \(= 0\)):
        \(\sqrt{\$500{,}000^2 + \$500{,}000^2} \approx \$707{,}000\)
  \item If perfectly negatively correlated:
        \(\$500{,}000 - \$500{,}000 = \$0\)
\end{itemize}

This is the principle of \textbf{diversification}\index{diversification}: combining uncorrelated or
negatively correlated risks reduces aggregate volatility below the sum of
individual volatilities. Insurers rely on this principle---pooling many
independent risks allows them to predict aggregate losses more reliably than
individual losses. However, correlations are difficult to estimate accurately
and may change during stress periods. As the LTCM case illustrates, risks that
appear uncorrelated in normal times may become highly correlated during
crises---and that is precisely when the stakes are highest.

\subsection{Probability Distributions and Loss Distributions}
\label{subsec:loss-distributions}

A \textbf{probability distribution}\index{probability distribution} describes the likelihood of different
possible outcomes. For risk management, loss distributions\index{loss distribution} describe the
probabilities of different loss amounts. Loss distributions can be
\emph{discrete}---taking on specific distinct values, like the number of
workplace accidents in a year (0, 1, 2, 3, \ldots)---or
\emph{continuous}---taking any value within a range, like the dollar amount of
property damage or a liability settlement.

The \textbf{shape} of the loss distribution is critical for risk management.
\emph{Symmetric distributions} are balanced around the mean, with equal
probability of values above and below the mean by similar amounts. The normal
distribution (bell curve) is the most familiar example. \emph{Right-skewed
(positively skewed) distributions} have an asymmetric long tail extending toward
larger values. Most losses are relatively small, but occasional large losses
create that long right tail. The mean exceeds the median because rare large
values pull the mean upward.

Right-skewed distributions characterize many insurable risks: most workplace
accidents result in minor injuries with modest costs, but occasional serious
injuries create very large claims. Most property damage incidents are minor, but
occasional fires or natural disasters create catastrophic losses. Understanding
whether a loss distribution is symmetric or skewed is essential for risk
management. If you assume losses are normally distributed when they are actually
right-skewed, you will underestimate the probability and magnitude of large
losses---a dangerous mistake, and one that contributed materially to LTCM's
failure.

% -----------------------------------------------------------------------
\section{Case Study: Quantifying Risk at Midtown Outfitters (Retailer)}
\label{sec:midtown-outfitters}
% -----------------------------------------------------------------------

To make these statistical concepts concrete, we work through a detailed case
of risk quantification for a regional retailer. This case demonstrates how an
organization gathers data, computes basic risk metrics, and interprets the
results to inform risk management decisions.

\subsection{Company Background and Risk Landscape}
\label{subsec:midtown-background}

Midtown Outfitters\index{Midtown Outfitters} is a regional chain of 45 retail clothing stores located
across three states in the Midwest. The company operates mall-based stores,
stand-alone locations, and an e-commerce platform. Annual revenues are
approximately \$180~million with 1,200 employees (800 full-time, 400
part-time). Like all retailers, Midtown Outfitters faces multiple risk
categories identified using the four-quadrant framework from Chapter~3.

\textbf{Hazard risks} include property damage to stores, inventory, and
fixtures from fire, weather, water damage, or vandalism; general liability from
customer slip-and-fall accidents or product liability claims; employee workplace
injuries; and vehicle accidents involving company-owned delivery vehicles.

\textbf{Operational risks} include inventory shrinkage from shoplifting,
employee theft, or administrative errors; supply chain disruptions affecting
product availability; point-of-sale system failures; and employee turnover and
training costs.

\textbf{Financial risks} include credit card processing fees and chargebacks;
currency exposure from imported merchandise; and interest rate risk on
variable-rate debt.

\textbf{Strategic risks} include e-commerce competition from larger national
retailers, changing consumer preferences, lease costs and commitments for mall
locations, and seasonal demand variability.

For this case study, we focus on quantifying several hazard and operational
risks where Midtown Outfitters has historical loss data: customer slip-and-fall
incidents, inventory shrinkage, and minor property damage.

\subsection{Data Collection and Organization}
\label{subsec:midtown-data}

Midtown Outfitters' risk manager begins by gathering internal data from multiple
sources. Incident reports from a centralized database record all reported
incidents---customer accidents, employee injuries, property damage events, and
security incidents---with date, location, description, immediate costs, and
whether a claim was filed. The company's insurance broker provides detailed
loss runs\index{loss runs} (historical claims data) for general liability, property, workers'
compensation, and other coverages. Loss runs include claim date, description,
status (open or closed), paid amounts, and reserves for open claims. Accounting
systems track inventory purchases, sales, and physical inventory counts, from
which shrinkage is calculated. Industry associations and insurance providers
publish benchmark statistics for retail loss frequencies and severities,
providing context for evaluating the company's experience relative to peers.
The risk manager organizes five years of data to ensure a reasonably large
sample while remaining recent enough to reflect current operations.

\subsection{Quantifying Slip-and-Fall Incidents}
\label{subsec:midtown-slipfall}

Customer slip-and-fall incidents represent Midtown Outfitters' most frequent
general liability exposure. The risk manager follows four analytical steps.

\textbf{Step~1: Calculate Frequency.} Over five years, Midtown Outfitters
experienced 68 slip-and-fall incidents across all stores: 12 incidents in
Year~1 (40 stores operating), 15 in Year~2 (42 stores), 14 in Year~3
(44 stores), 13 in Year~4 (45 stores), and 14 in Year~5 (45 stores). Total
exposure is 215 store-years (40\,+\,42\,+\,44\,+\,45\,+\,45).

\[
  \text{Frequency} = \frac{68 \text{ incidents}}{215 \text{ store-years}}
  = 0.316 \text{ incidents per store per year}
\]

Each store experiences roughly one incident every three years. For the current
45 stores, expected annual incidents \(= 45 \times 0.316 = 14.2\).

\textbf{Step~2: Analyze Severity Distribution.} Of the 68 incidents: 52
incidents (76\%) resulted in no claim, with average cost \$300 per incident
for administrative time and minor accommodation; 12 incidents (18\%) resulted
in small claims averaging \$4,500 each; and 4 incidents (6\%) resulted in
larger claims costing \$15,000, \$22,000, \$35,000, and \$48,000 respectively.

\begin{align*}
  \text{Total costs} &= (52 \times \$300) + (12 \times \$4{,}500)
                       + (\$15{,}000 + \$22{,}000 + \$35{,}000 + \$48{,}000)\\
                     &= \$15{,}600 + \$54{,}000 + \$120{,}000 = \$189{,}600
\end{align*}

Average severity \(= \$189{,}600 / 68 = \$2{,}788\) per incident. Average
severity is pulled up by the four large claims even though most incidents
(76\%) cost only \$300. The median severity is \$300, illustrating the
right-skewed nature of the loss distribution.

\textbf{Step~3: Calculate Expected Loss.}
\[
  \text{Expected annual loss} = 0.316 \times \$2{,}788 = \$881 \text{ per store per year}
\]
For 45 stores: \(45 \times \$881 = \$39{,}645\) expected annual loss from
slip-and-fall incidents.

\textbf{Step~4: Assess Variability.} Historical annual total costs varied
considerably: Year~1 \$35,000 (including one \$22,000 claim), Year~2 \$42,000,
Year~3 \$18,000 (mostly minor incidents), Year~4 \$62,000 (including a \$35,000
claim), and Year~5 \$32,600. Mean annual loss is \$37,920 with standard
deviation approximately \$16,000. The high standard deviation relative to the
mean reflects significant year-to-year variability driven by whether any large
claims occur in a given year.

\subsection{Quantifying Inventory Shrinkage}
\label{subsec:midtown-shrinkage}

Inventory shrinkage\index{inventory shrinkage}---the difference between book inventory and physical
inventory---represents ongoing operational loss from shoplifting, employee
theft, vendor fraud, and administrative errors. Midtown Outfitters conducts
annual physical inventory counts at each store. Historical shrinkage as a
percentage of sales was: Year~1 2.1\%, Year~2 1.8\%, Year~3 2.3\%, Year~4
1.9\%, and Year~5 2.0\%, yielding an average rate of 2.02\%. With annual sales
of \$180~million, expected annual shrinkage is \(\$180\text{M} \times 2.02\% =
\$3.64\text{M}\).

Shrinkage is relatively predictable---the standard deviation of the five annual
percentages is only 0.19 percentage points---because it accumulates through
many small incidents rather than clustering in rare large events. Unlike
slip-and-fall claims, where a single large settlement can dominate a given year,
shrinkage aggregates gradually. That said, shrinkage varies considerably across
stores, and high-shrinkage locations can be identified for enhanced security,
employee training, or store design modifications.

\subsection{Quantifying Property Damage}
\label{subsec:midtown-property}

Minor property damage---vandalism, broken windows, roof leaks, fixture
damage---occurs periodically. Major losses are rare but potentially
catastrophic, creating a right-skewed distribution. Over five years, Midtown
Outfitters experienced 85 minor incidents (average cost \$1,200 each), 8
moderate incidents (average cost \$8,500 each), and 1 major incident---a store
damaged by a tornado---costing \$180,000. Total incidents were 94 over 215
store-years.

\begin{align*}
  \text{Frequency} &= 94 / 215 = 0.437 \text{ per store per year}\\
  \text{Total costs} &= (85 \times \$1{,}200) + (8 \times \$8{,}500) + \$180{,}000\\
                     &= \$102{,}000 + \$68{,}000 + \$180{,}000 = \$350{,}000\\
  \text{Average severity} &= \$350{,}000 / 94 = \$3{,}723\\
  \text{Expected annual loss} &= 0.437 \times \$3{,}723 = \$1{,}627 \text{ per store}
                                 = \$73{,}215 \text{ for 45 stores}
\end{align*}

This distribution is heavily right-skewed. Excluding the tornado, the 93 other
incidents averaged only \$1,828 in severity. The single tornado caused more
damage than all other incidents combined over five years. Most of the time,
property losses are minor and manageable; rare severe events can dominate total
losses over any multi-year period.

\subsection{Risk Matrix and Prioritization}
\label{subsec:midtown-matrix}

To integrate quantitative analysis with qualitative assessment, Midtown
Outfitters creates a \(3 \times 3\) risk matrix\index{risk matrix} plotting risks based on
likelihood (frequency) and impact (severity). Likelihood categories are high
(more than 5 incidents per year across all stores), medium (1--5 per year), and
low (fewer than 1 per year). Impact categories are high (potential single-event
loss exceeding \$100,000), medium (\$10,000--\$100,000), and low (typically
under \$10,000).

Based on the quantitative analysis: slip-and-fall incidents fall into the
Medium-Medium cell (14 expected incidents per year, most claims under \$50,000
but potential for larger); inventory shrinkage falls into the High-Low cell
(occurs continuously, each incident small, but aggregate annual impact of
\$3.64~million demands attention); minor and moderate property damage falls in
the High-Low cell; and catastrophic property events fall in the Low-High cell
(a tornado once in five years, costing \$180,000 or more). Risks in
Medium-Medium and Low-High cells warrant particular attention. Medium-Medium
risks accumulate substantial aggregate losses. Low-High risks, while rare, can
create severe financial stress if uninsured.

\subsection{Management Interpretation and Decisions}
\label{subsec:midtown-decisions}

The quantitative analysis informs several management decisions. For
\textbf{insurance purchasing}, expected annual losses from slip-and-fall
(\$39,645) and typical property damage (\$73,215) are modest relative to
revenues, suggesting these risks can be largely retained through deductibles.
However, catastrophic losses and large liability claims justify insurance for
severe events. Management decides to maintain general liability insurance with a
\$25,000 per-occurrence retention and property insurance with a \$10,000
per-occurrence deductible---retaining most frequent losses while transferring
tail risk.

For \textbf{loss prevention}, the \$3.64~million annual shrinkage represents
2\% of sales and a meaningful profit impact. Even modest reduction---from
2.02\% to 1.75\%---would save approximately \$500,000 annually, justifying
substantial investment in electronic article surveillance, employee training,
and store design modifications.

For \textbf{slip-and-fall prevention}, analysis shows frequency spikes during
winter months (ice and snow tracked into stores) and on rainy days (wet
floors). Management implements protocols for prompt floor drying, warning signs,
and enhanced entrance matting during adverse weather. Management also
establishes quarterly monitoring of shrinkage rates, incident frequencies, and
claims severity; deviations from expected levels trigger investigation. This
case illustrates the cycle of risk quantification: gather data, calculate
frequency and severity, compute expected losses and variability, use risk
matrices to prioritize, and translate analysis into management decisions.

% -----------------------------------------------------------------------
\section{Case Study: Quantifying Risk at Precision Parts Inc.\ (Manufacturer)}
\label{sec:precision-parts}
% -----------------------------------------------------------------------

Manufacturing companies face different risk profiles than retailers, with
emphasis on equipment reliability, product quality, workplace safety, and
supply chain continuity. This second case demonstrates risk quantification in a
manufacturing context.

\subsection{Company Background}
\label{subsec:precision-background}

Precision Parts Inc.\index{Precision Parts Inc.}\ manufactures metal components for the automotive and
aerospace industries. The company operates two production facilities (Michigan
and Ohio) employing 650 workers. Annual revenues are approximately
\$120~million. Production involves machining, stamping, welding, and finishing
operations using specialized equipment. Key risks include equipment breakdown
and unplanned production downtime; product defects leading to customer returns,
warranty claims, or recalls; workplace injuries in production environments; and
supply chain disruptions for raw materials (primarily specialty steel and
aluminum).

\subsection{Quantifying Equipment Breakdown and Downtime}
\label{subsec:precision-equipment}

Equipment reliability directly affects production capacity, customer deliveries,
and profitability. Precision Parts maintains detailed maintenance logs for all
equipment failures, repair times, and costs. Analysis of three years of data
for critical production equipment (CNC machines, stamping presses, welding
robots) shows 72 total failures (24, 21, and 27 in years 1--3), averaging 24
per year. Failure categories are minor failures (48 of 72, 67\%), averaging
4~hours of downtime and repaired by maintenance staff; moderate failures (18 of
72, 25\%), averaging 16~hours and requiring replacement parts; and major
failures (6 of 72, 8\%), averaging 72~hours and requiring extensive repair or
equipment replacement.

\[
  \text{Average downtime per failure}
  = \frac{(48 \times 4) + (18 \times 16) + (6 \times 72)}{72}
  = \frac{192 + 288 + 432}{72}
  = 12.67 \text{ hours}
\]

Expected annual downtime \(= 24 \times 12.67 = 304\) hours. Downtime costs
include lost contribution margin, idled labor, rush charges, and potential late
delivery penalties. Precision Parts estimates \$1,200 per hour of downtime,
yielding expected annual downtime cost of \(304 \times \$1{,}200 = \$364{,}800\).
The timing of downtime matters, however: downtime during peak demand periods
costs more than downtime during slack periods, and downtime affecting equipment
with no backup is more costly than downtime on equipment with redundancy.

\subsection{Quantifying Product Defects and Warranty Claims}
\label{subsec:precision-defects}

Product quality failures create direct costs (scrap, rework, customer returns)
and potential liability if defective parts cause customer production problems or
safety issues. Precision Parts tracks defect rates at three stages: internal
detection before shipment, customer returns after shipment, and warranty claims
after customer installation.

Over three years and 14.5~million parts produced, quality control rejected
87,000 units (0.60\% internal defect rate), yielding scrap cost of
\(0.60\% \times \$120\text{M} = \$720{,}000\) per year. Of 14.413~million
parts shipped, 5,760 were returned by customers (0.04\% return rate), with
average returns of 1,920 per year at \$180 per return (\$345,600 per year).
Of 84 warranty claims over three years (28 per year average): 64 claims (76\%)
were minor part replacements averaging \$2,500; 16 claims (19\%) were moderate
failures causing customer downtime averaging \$12,000; and 4 claims (5\%) were
major failures causing customer production disruption averaging \$45,000.

\[
  \text{Annual warranty cost}
  = \frac{(64 \times \$2{,}500) + (16 \times \$12{,}000) + (4 \times \$45{,}000)}{3}
  = \frac{\$532{,}000}{3}
  = \$177{,}333
\]

Total quality-related costs sum to approximately \$1.24~million per year
(about 1.0\% of sales). The four major warranty claims at \$45,000 each
represent only 5\% of claim frequency but 34\% of total warranty costs---a
right-skewed distribution mirroring the pattern observed in slip-and-fall claims
and property damage.

\subsection{Quantifying Workplace Injury Risk}
\label{subsec:precision-injuries}

Over three years, Precision Parts recorded 48 recordable workplace injuries
from 3.9 million total work hours. The OSHA\index{OSHA} injury rate is
\((48 / 3{,}900{,}000) \times 200{,}000 = 2.46\) per 200,000 work hours,
favorably below the manufacturing industry average of approximately 3.5,
suggesting Precision Parts' safety program is effective. Injury severity
distribution shows 32 minor injuries (67\%) averaging \$2,800 each, 12
moderate injuries (25\%) averaging \$18,000 each, and 4 serious injuries (8\%)
averaging \$62,000 each. Average annual direct injury cost is
\(\$553{,}600 / 3 = \$184{,}533\). Indirect costs---lost productivity, training
replacements, incident investigation, potential OSHA penalties---approximately
double total costs, suggesting true annual workplace injury cost around
\$370,000.

\subsection{Risk Correlations and Compounding Effects}
\label{subsec:precision-correlations}

The risk manager recognizes that manufacturing risks can be correlated\index{correlation}. During
periods of rushed production to meet delivery deadlines, equipment operates at
higher speeds and longer hours (increasing breakdown frequency), quality
pressure increases error rates, and overtime fatigue increases accident rates.
Statistical analysis reveals positive correlation between equipment breakdowns
and defect rates (correlation coefficient approximately \(+0.35\)), suggesting
common causes---rushed production and inadequate maintenance---affect both
simultaneously.

Simply adding the individual risk measures (\$364,800 + \$1,242,933 + \$184,533
\(= \$1{,}792{,}266\)) understates total risk because it ignores that bad
outcomes tend to cluster. During stress periods, equipment failures cause missed
deliveries, pressure to catch up increases defect rates, and rushed work
increases injuries. This insight leads to process improvements targeting root
causes: preventive maintenance schedules that reduce rush periods, capacity
buffers that allow normal-paced production even when demand peaks, and quality
gates that halt production if defect rates exceed thresholds.

\subsection{Qualitative Risk Assessment for Manufacturing}
\label{subsec:precision-qualitative}

Beyond quantitative analysis, Precision Parts uses qualitative tools to assess
risks that are difficult to quantify. \textbf{Scenario analysis}\index{scenario analysis} for supply
chain disruption has managers develop scenarios for supply interruptions
(supplier bankruptcy, natural disaster at a supplier facility, trade
restrictions) and estimate potential duration, workarounds available, customer
impacts, and revenue loss. The exercise reveals that critical steel supplier
concentration creates substantial risk: loss of this supplier could halt
production for weeks. \textbf{Expert elicitation}\index{expert elicitation} for catastrophic equipment
failure has maintenance engineers estimate a 2--5\% annual probability of
complete failure of the firm's custom stamping press, with potential impact of
\$2--3~million in lost contribution margin during extended downtime. This
qualitative assessment justifies maintaining spare critical components despite
high inventory costs. A \textbf{risk matrix} displays both quantitative and
qualitative assessments visually, helping senior management prioritize attention
across diverse risks. Precision Parts' experience shows that effective risk
assessment combines quantitative analysis where data exists with qualitative
judgment for low-frequency, high-consequence events and emerging risks that
lack historical data.

% -----------------------------------------------------------------------
\section{Loss Distributions and Skewness}
\label{sec:loss-distributions}
% -----------------------------------------------------------------------

Understanding the shape of loss distributions---not just the mean and standard
deviation---is crucial for risk management. This section explores normal versus
skewed distributions and explains why skewness matters.

\subsection{The Normal Distribution}
\label{subsec:normal-distribution}

The normal distribution\index{normal distribution} (bell curve) is symmetric around its mean. Values are
equally likely to fall above or below the mean by similar amounts, and the
distribution is characterized by mean = median = mode; 68~percent of values
within one standard deviation of the mean; 95~percent within two standard
deviations; and 99.7~percent within three standard deviations. Many natural
phenomena follow approximately normal distributions due to the central limit
theorem\index{central limit theorem}: when many independent small factors combine to determine an outcome,
the distribution tends toward normal. For risk management, normal distributions
are relatively ``well-behaved.'' Extreme outcomes are rare and become
exponentially less likely as you move further from the mean. A normal
distribution with mean \$100,000 and standard deviation \$20,000 has only
0.13~percent probability of exceeding \$160,000 (three standard deviations).

\subsection{Right-Skewed Distributions and Tail Risk}
\label{subsec:right-skewed}

Many insurable losses follow right-skewed distributions\index{right-skewed distribution}. Most loss events are
small, but occasional large losses create a long right tail. Key characteristics
include: mean exceeds median (rare large values pull mean upward); substantial
probability mass in the right tail; and the maximum observed loss does
\emph{not} bound future possibilities.

Consider the property damage distribution from Midtown Outfitters in
Section~\ref{subsec:midtown-property}: 85 minor incidents averaging \$1,200,
8 moderate incidents averaging \$8,500, and 1 catastrophic incident costing
\$180,000. The median severity is \$1,200, but mean severity is \$3,723. The
single tornado caused more than half of all property damage costs over five
years. Right-skewed distributions are common for property damage (most events
minor; occasional fires or catastrophes), liability claims (most claims modest;
occasional million-dollar judgments), product recalls (most quality issues
minor; occasional large-scale recalls), and operational losses (most errors
immaterial; occasional major frauds or systems failures).

\subsection{Why Skewness Matters for Risk Management}
\label{subsec:skewness-matters}

\textbf{Underestimating large loss potential.} If you assume losses follow a
normal distribution when they are actually right-skewed, you will severely
underestimate the probability of large losses. This is precisely what happened
at LTCM: models calibrated to historical volatility treated extreme events as
nearly impossible when they were, in fact, quite possible. The ``six-sigma''
events that LTCM's models dismissed as absurdly unlikely occurred because the
true distribution had fatter tails than the historical data suggested.

\textbf{Mean understates typical experience.} In right-skewed distributions,
the mean exceeds the median, so ``typical'' years experience losses below the
mean. The mean is pulled upward by occasional years with very large losses.
Budgeting based on mean losses may provide inadequate reserves in bad years.

\textbf{Tail events dominate total losses.}\index{tail risk} For highly skewed distributions, a
small percentage of events may account for most total losses. In Midtown
Outfitters' property damage experience, the single tornado (1\% of incidents)
caused 51\% of total losses. Risk management must focus on preventing or
mitigating tail events, not just managing typical losses.

\textbf{Insurance value is concentrated in tail protection.} If an organization
can absorb typical losses but not catastrophic losses, insurance provides
greatest value by covering the tail. A deductible that retains small and
moderate losses while transferring large losses aligns with this logic:
self-insure the predictable part, transfer the catastrophic risk.

\subsection{Measuring and Describing Skewness}
\label{subsec:measuring-skewness}

Statisticians measure skewness\index{skewness} formally, but qualitative assessment often
suffices for risk management. \emph{Positive skewness} (right-skewed) is
characterized by mean exceeding median and a long right tail; it is
characteristic of most insurable losses. \emph{Negative skewness} is
characterized by mean below median and a long left tail; it is less common for
losses but occurs for some financial returns. \emph{Zero skewness} (symmetric)
has mean equal to median; the normal distribution has zero skewness. Beyond
skewness, \textbf{kurtosis}\index{kurtosis} measures the fatness of tails. Distributions with
high kurtosis (leptokurtic distributions) have fatter tails than the normal
distribution---extreme events occur more frequently than a normal distribution
would predict. Many financial market returns exhibit fat tails\index{fat tails}, meaning crashes
and booms occur more often than the normal model suggests. LTCM's losses
illustrated this at an extraordinary scale.

% -----------------------------------------------------------------------
\section{Excel-Based Workers' Compensation Case: Estimating Value at Risk}
\label{sec:northland-var}
% -----------------------------------------------------------------------

This section presents a detailed teaching case suitable for Excel analysis.
The case demonstrates how to build an empirical loss distribution from
historical data and estimate Value at Risk (VaR)---a widely used summary risk
measure.

\subsection{Company Background: Northland Manufacturing}
\label{subsec:northland-background}

Northland Manufacturing\index{Northland Manufacturing} operates a metal fabrication facility employing
400 workers in production, maintenance, and administrative roles. The company
has maintained detailed workplace injury records for ten years as required by
OSHA. Management wants to quantify workers' compensation\index{workers' compensation} risk to inform
insurance purchasing decisions, retained loss budgeting, and safety program
prioritization.

\subsection{Data Description}
\label{subsec:northland-data}

Northland experienced 138 recordable workplace injuries over ten years. For
each injury, data includes year of occurrence, lost workdays (if any), medical
costs, indemnity costs (wage replacement for lost time), and total claim cost.
Table~\ref{tab:northland-annual} summarizes annual experience.

\begin{ermtable}{Annual Workers' Compensation Experience, Northland Manufacturing}[tab:northland-annual]
  \begin{tabularx}{\linewidth}{cYYcc}
    \toprule
    Year & Employees & Work Hours & Injuries & Total Cost \\
    \midrule
    1  & 380 & 760,000 & 18 & \$187,500 \\
    2  & 390 & 780,000 & 12 & \$142,800 \\
    3  & 395 & 790,000 & 15 & \$256,400 \\
    4  & 400 & 800,000 & 11 & \$128,900 \\
    5  & 400 & 800,000 & 16 & \$198,600 \\
    6  & 410 & 820,000 & 14 & \$172,300 \\
    7  & 405 & 810,000 & 13 & \$283,700 \\
    8  & 400 & 800,000 & 10 & \$115,800 \\
    9  & 395 & 790,000 & 15 & \$224,500 \\
    10 & 400 & 800,000 & 14 & \$189,200 \\
    \midrule
    \textbf{Total} & \textit{3{,}975 avg.} & 7{,}950{,}000 & 138 & \$1{,}899{,}700 \\
    \bottomrule
  \end{tabularx}
  \tablenote{Data are hypothetical and intended for instructional purposes.}
\end{ermtable}

The individual claim distribution shows 98 claims (71\%) as minor injuries,
medical-only with no lost time, ranging from \$800 to \$5,000 with mean
\$2,400; 32 claims (23\%) as moderate injuries with some lost time, ranging
from \$6,000 to \$20,000 with mean \$11,500; and 8 claims (6\%) as serious
injuries with extended lost time, ranging from \$25,000 to \$75,000 with mean
\$42,000.

\subsection{Step-by-Step Excel Analysis}
\label{subsec:northland-excel}

\textbf{Step~1: Calculate Frequency Rate.}

\[
  \text{Rate per 200{,}000 hours}
  = \frac{138}{7{,}950{,}000} \times 200{,}000 = 3.47 \text{ injuries per 200,000 hours}
\]

Per employee per year: \(138 / (3{,}975 \times 10) \approx 0.00347\), or
3.47 per 1,000 employees per year. For the current workforce of 400 employees,
expected annual injuries \(= 400 \times 0.00347 \approx 14\), consistent with
the historical average.

\begin{excelblock}
=COUNT(injuries)/SUM(work\_hours)       \% frequency per hour\\
=Current\_employees * Frequency * 2000  \% expected annual injuries
\end{excelblock}

\textbf{Step~2: Build Severity Distribution.} Average severity across all 138
claims is \(\$1{,}899{,}700 / 138 = \$13{,}765\). Median severity is
approximately \$2,800 (71\% of claims fall below \$5,000). The mean
(\$13,765) exceeds the median (\$2,800) by nearly 5$\times$, confirming
right-skewed distribution.

\begin{excelblock}
=AVERAGE(claim\_amounts)  \% mean severity\\
=MEDIAN(claim\_amounts)   \% median severity
\end{excelblock}

\textbf{Step~3: Estimate Expected Annual Loss.}

\[
  \text{Expected annual loss} = 14 \text{ injuries} \times \$13{,}765
  = \$192{,}710
\]

Historical annual losses ranged from \$115,800 (Year~8) to \$283,700 (Year~7),
with mean \$189,970 and standard deviation approximately \$53,000. The expected
loss of \$192,710 aligns closely with the historical mean, as expected.

\textbf{Step~4: Simulate Annual Loss Distribution.} Two approaches exist.
The \emph{historical simulation approach}\index{historical simulation} uses the historical annual losses
directly. Sorted from smallest to largest: \$115,800; \$128,900; \$142,800;
\$172,300; \$187,500; \$189,200; \$198,600; \$224,500; \$256,400; \$283,700.
The \emph{Monte Carlo simulation approach}\index{Monte Carlo simulation} simulates many possible years by
randomly drawing the number of injuries (using a Poisson distribution\index{Poisson distribution} calibrated
to historical frequency) and then randomly drawing a severity for each injury
from the historical severity distribution, repeating 10,000 times to build a
full distribution of possible annual losses.

\begin{excelblock}
Column A: Simulation number (1 to 10,000)\\
Column B: Random injuries: =POISSON.DIST(14,RAND(),FALSE)\\
Column C: For each injury, random severity via lookup table\\
Column D: Total annual loss: =SUM(injury severities for that year)\\
\medskip
90th percentile: =PERCENTILE(Column\_D, 0.90)\\
95th percentile: =PERCENTILE(Column\_D, 0.95)
\end{excelblock}

\subsection{Calculating Value at Risk}
\label{subsec:northland-var-calc}

\textbf{Value at Risk (VaR)}\index{Value at Risk (VaR)} answers the question: how bad could losses be in
a bad year? Specifically, VaR at a given confidence level (say, 95\%) is the
threshold that losses will exceed with only 5\% probability.

\textbf{VaR at 90\% confidence:} The loss level exceeded in only 10\% of
years. From the sorted historical data, approximately \$256,400 (the ninth
largest loss in ten years).

\begin{excelblock}
=PERCENTILE.INC(annual\_losses, 0.90)  \% returns approximately \$256,400
\end{excelblock}

\emph{Interpretation:} With 90\% confidence, annual workers' compensation
losses will not exceed \$256,400, assuming future resembles past.

\textbf{VaR at 95\% confidence:} With only ten data points, interpolation is
required.

\begin{excelblock}
=PERCENTILE.INC(annual\_losses, 0.95)  \% returns approximately \$270,000
\end{excelblock}

\textbf{VaR at 99\% confidence:} With only ten years of data, we have not
observed a 1-in-100-year event, so estimation requires distributional
assumptions. Using a fitted lognormal\index{lognormal distribution} or gamma distribution calibrated to our
data, we might estimate the 99\% VaR at \$400,000--\$500,000.

\subsection{Interpreting and Using VaR}
\label{subsec:northland-interpretation}

\textbf{Risk budgeting and capital planning.} To be 95\% confident of having
sufficient reserves, Northland should maintain approximately \$270,000 in
reserves or insurance coverage beyond expected losses. This does not mean the
company will spend \$270,000---in most years it will spend considerably less.
It means setting aside enough to handle a bad-but-not-catastrophic year.

\textbf{Insurance deductible selection.} If Northland retains losses up to
\$150,000 (below the 90\% VaR), it will be able to cover losses from its own
resources in more than 90\% of years. Losses exceeding \$150,000---occurring
perhaps once every 6--10 years---would be covered by insurance. This structure
balances the cost of retaining predictable frequency against the protection
value of transferring tail severity.

\textbf{Comparative risk assessment.} VaR enables comparing different risks on
a common basis. If workplace injury VaR (95\%) is \$270,000 while property
damage VaR (95\%) is \$150,000, workplace injury represents the greater tail
risk and may warrant more management attention and resources.

% -----------------------------------------------------------------------
\section{Value at Risk: Usefulness and Limitations}
\label{sec:var-limitations}
% -----------------------------------------------------------------------

VaR has become the most widely used summary risk measure in finance and
increasingly in enterprise risk management. Understanding both its value and
its limitations is essential for proper use.

\subsection{Why VaR Is Appealing}
\label{subsec:var-appeal}

\textbf{Simplicity and clarity.} VaR condenses a complex distribution into a
single number that is easy to communicate: ``Losses will not exceed \$X with
probability Y\%.'' This simplicity makes VaR understandable to senior
management and boards without statistical training.

\textbf{Common risk language.} VaR provides a standardized metric for
comparing different risks across the enterprise, over time, and across
companies. A bank can compare market risk VaR, credit risk VaR, and operational
risk VaR on common terms. A non-financial firm can compare workers' compensation
VaR against property VaR using the same methodology.

\textbf{Forward-looking downside focus.} Unlike average losses (which describe
central tendency), VaR focuses on adverse outcomes---how bad things might be.
This aligns with risk management's purpose of preparing for adverse scenarios
rather than simply describing historical experience.

\textbf{Regulatory acceptance.} Basel banking regulations require VaR
calculations for market risk. Insurance solvency regulations use similar tail
risk measures. Regulatory acceptance has driven widespread adoption and
standardization across industries.

\subsection{Limitations of VaR}
\label{subsec:var-limits}

Despite these advantages, VaR has significant limitations\index{Value at Risk (VaR)!limitations}. The LTCM story in
the opening case illustrates what can go wrong.

\textbf{VaR says nothing about losses beyond the threshold.} VaR at 95\%
confidence tells you the threshold exceeded 5\% of the time, but it does not
tell you \emph{how much} losses might exceed that threshold. Consider two
portfolios: Portfolio~A has 95\% VaR of \$1~million and losses ranging from
\$1~million to \$1.2~million in the worst 5\% of years; Portfolio~B has the
same 95\% VaR but losses ranging from \$1~million to \$10~million in the worst
5\% of years. Both portfolios have identical VaR. Portfolio~B is far riskier.
VaR alone does not distinguish these cases.

\textbf{Model and assumption dependence.} VaR estimates depend heavily on
which model and assumptions are used. Historical simulation VaR assumes the
future resembles the past. Parametric VaR assumes losses follow a particular
distribution. Different methods can yield substantially different estimates for
the same underlying risk. LTCM's models assumed stable correlations; when
correlations spiked in 1998, the models were badly wrong and the firm was
catastrophically leveraged against that error.

\textbf{Data limitations.} VaR estimates are unreliable when historical data is
limited, particularly when estimating high-confidence VaR (99\% or higher) that
requires predicting rare events. With only ten years of data, estimating 99\%
VaR requires extreme extrapolation.

\textbf{Stress correlations.} VaR models typically estimate correlations from
normal market conditions. During crises, correlations increase as multiple
problems emerge simultaneously. VaR models may systematically underestimate
aggregate risk during stress periods---precisely when accurate risk measurement
matters most.

\textbf{Incentive effects.} Because VaR is a threshold measure, traders or
managers might structure positions that minimize VaR while taking substantial
risk just beyond the threshold---satisfying risk limits while actually
increasing danger.

\subsection{Complementary Risk Measures}
\label{subsec:complementary-measures}

\textbf{Expected Shortfall (Conditional VaR or CVaR)}\index{Expected Shortfall (ES)} is the average loss given
that losses exceed VaR. It answers: if we have a bad year (losses exceeding
VaR), how bad will it be on average? For Northland Manufacturing, VaR (90\%) is
\$256,400. Expected Shortfall (90\%) = average of losses exceeding \$256,400 =
\((\$256{,}400 + \$283{,}700) / 2 = \$270{,}050\). When annual losses exceed
the 90\% threshold, they average \$270,000---not drastically more severe than
the threshold, suggesting the tail is not extremely fat for this dataset.

\textbf{Stress testing and scenario analysis}\index{stress testing}\index{scenario analysis} develop specific adverse scenarios
(``What if a major fire destroys our primary facility?'' or ``What if our
largest customer declares bankruptcy?'') and assess impacts. Scenarios
complement VaR by exploring specific threats that might not appear likely based
on historical data but represent genuine vulnerabilities.

\textbf{Sensitivity analysis}\index{sensitivity analysis} tests how risk measures change if key assumptions
change. If VaR estimates are highly sensitive to assumed distribution or
correlation assumptions, this reveals model uncertainty requiring caution in
relying on the point estimate.

\subsection{Integrating VaR into Risk Management}
\label{subsec:var-integration}

\textbf{Use VaR for relative comparisons and trends.} Track how VaR changes
over time as risk exposure evolves. Compare VaR across different business units
or risk types. VaR is more reliable for these comparative purposes than for
precise absolute risk quantification.

\textbf{Supplement with stress testing.} Develop worst-case scenarios beyond
VaR thresholds. If VaR (99\%) is \$500,000, ask: what is our maximum possible
loss if everything goes wrong simultaneously? This is the question LTCM's
management failed to take seriously enough.

\textbf{Back-test VaR estimates.}\index{back-testing} Compare actual losses to VaR predictions. If
VaR (95\%) should be exceeded 5\% of the time but is actually exceeded 15\% of
the time, the model is mis-calibrated and requires adjustment.

\textbf{Maintain qualitative judgment.} VaR should inform decisions, not make
them automatically. Quantitative measures discipline thinking but do not
replace it. A model that says the loss threshold is \$270,000 does not tell you
what to do about it---that requires judgment about the firm's financial
capacity, strategic priorities, and risk appetite.

% -----------------------------------------------------------------------
\section{Limits of Quantitative Measures and the Role of Qualitative Assessment}
\label{sec:qual-quant-integration}
% -----------------------------------------------------------------------

\subsection{When Quantification Is Most Valuable}
\label{subsec:when-quantification}

Quantification works best when sufficient historical data exists---frequent,
relatively homogeneous events such as workplace accidents, warranty claims, and
minor property damage generate data amenable to statistical analysis; when the
future resembles the past---stable processes and environments mean historical
patterns will likely continue; when risks are measurable in common units so
losses can be expressed in dollars, enabling aggregation and comparison; and
when stakeholders expect quantification, as when regulators, boards, or
investors require numerical risk assessments.

\subsection{When Qualitative Assessment Is Essential}
\label{subsec:when-qualitative}

\textbf{Rare, unprecedented, or emerging risks.} Catastrophic equipment
failures, novel cybersecurity threats, strategic disruptions from new
competitors, or regulatory changes may have little or no historical precedent.
Quantifying these risks requires expert judgment, scenario analysis, and
structured discussion rather than statistical modeling.

\textbf{Complex, interconnected risks.} When multiple risks interact in ways
that are difficult to model quantitatively, scenario workshops where diverse
experts explore ``what if'' questions can reveal vulnerabilities that
statistical models miss.

\textbf{Reputational and intangible risks.} Damage to brand reputation, loss
of customer trust, or erosion of company culture create real consequences but
resist precise measurement. Qualitative assessment through stakeholder
interviews, sentiment analysis, and narrative scenario development provides
insight that numbers alone cannot.

\textbf{Black swan events.}\index{black swan events} Extremely rare, catastrophic events that have never
occurred---or occurred so rarely that statistical analysis provides little
guidance---require qualitative scenario thinking rather than extrapolation from
historical data. Climate change creating unprecedented weather patterns,
pandemics, geopolitical shocks, and financial market crashes all fall in this
category. LTCM's management did not adequately ask what would happen if
liquidity evaporated across all markets simultaneously---a scenario without
close historical precedent.

\subsection{Integrating Quantitative and Qualitative Assessment}
\label{subsec:qq-integration}

The most effective risk assessment uses both approaches systematically.
\textbf{Use quantitative analysis as foundation}: for risks with adequate data,
compute frequency, severity, expected loss, and variability. This provides
objective baseline assessment and enables tracking changes over time.
\textbf{Overlay qualitative judgment}: subject matter experts adjust
quantitative estimates based on factors not captured in historical data---changes
in operations, new controls, emerging threats, or risk correlations. Expert
judgment refines rather than replaces statistical analysis.
\textbf{Communicate with both numbers and narratives}: present risk assessment
to decision-makers using quantitative metrics and qualitative descriptions.
Different audiences respond differently, and narratives make the numbers
interpretable. \textbf{Use risk matrices and heat maps}\index{risk matrix}\index{heat map} to plot risks on axes
of likelihood and impact, enabling prioritization and a portfolio view of all
organizational risks. \textbf{Review and update regularly}: as new data
accumulates, update quantitative models; as the business environment changes,
update qualitative assessments.

\subsection{The Art and Science of Risk Assessment}
\label{subsec:art-science}

Risk assessment is both art and science. The science involves rigorous data
analysis, statistical methods, and systematic processes. The art involves
judgment about assumptions, interpretation of results, assessment of factors
not captured quantitatively, and communication that drives action. Effective
risk managers develop both capabilities: the technical skills to conduct
quantitative analysis properly and the judgment to recognize when quantification
misleads, when assumptions are questionable, and when qualitative insight should
override quantitative models. The Moneyball analogy is apt here: the Oakland
A's did not stop using scouts. They used quantitative metrics to discipline what
the scouts were looking for, and they maintained qualitative judgment about fit,
character, and factors their models could not capture. The best risk managers
do the same.

\begin{center}\rule{0.5\linewidth}{0.5pt}\end{center}

% -----------------------------------------------------------------------
\begin{keytakeaways}
\item
  \textbf{Risk Quantification vs.\ Risk Qualification}: Quantification uses
  data and statistics to measure risk objectively; qualification uses judgment,
  scenarios, and narrative assessment. Both are necessary; neither is
  sufficient alone.
\item
  \textbf{Frequency, Severity, and Expected Loss}: How often adverse events
  occur (frequency) and how costly each event is (severity) combine to
  determine expected loss. Standard deviation measures variability around that
  expectation; higher standard deviation means greater uncertainty.
\item
  \textbf{Quantification Applies Broadly}: The frequency-severity framework is
  not limited to financial institutions. Hospitals, airlines, manufacturers,
  and sports franchises have all improved resource allocation by measuring
  adverse events rigorously against a relevant exposure base and comparing
  against benchmarks.
\item
  \textbf{Correlations Shape Aggregate Risk}: Uncorrelated risks diversify,
  reducing aggregate volatility below the sum of individual volatilities.
  Correlated risks compound, especially in stress periods---as LTCM
  discovered. Historical correlations may not hold during crises.
\item
  \textbf{Skewed Distributions and Tail Risk}: Many insurable losses follow
  right-skewed distributions. Mean severity exceeds median severity; tail
  events can dominate total losses despite being rare. Assuming normality when
  distributions are skewed leads to systematic underestimation of catastrophic
  loss potential.
\item
  \textbf{Value at Risk (VaR)}: VaR estimates a threshold that losses will
  exceed with only small probability (e.g., 5\%). It is useful for capital
  planning, insurance purchasing, and stakeholder communication. But VaR
  ignores losses beyond the threshold, depends heavily on models and
  assumptions, and fails most conspicuously when it matters most.
\item
  \textbf{LTCM as a Cautionary Model}: Sophisticated quantification at LTCM
  transformed model risk into the central risk factor. The assumptions that
  made the VaR numbers reassuring---stable distributions, low cross-asset
  correlations, liquid exit markets---all failed simultaneously in 1998. Sizing
  positions based on optimistic model estimates under high leverage is
  effectively betting the firm on the model being right.
\item
  \textbf{Qualitative Judgment Remains Essential}: Rare events, emerging risks,
  and complex interdependencies resist quantification. Stress testing, scenario
  analysis, and expert elicitation address what historical data cannot. The
  goal is informed risk assessment, not mathematical precision for its own
  sake.
\end{keytakeaways}

% -----------------------------------------------------------------------
\begin{keyterms}
  \item[\keyterm{Expected loss}]
    The average loss anticipated over a specified period, calculated as
    frequency \(\times\) average severity.
  \item[\keyterm{Expected Shortfall (ES)}]
    The average loss given that losses exceed the VaR threshold; also called
    Conditional VaR (CVaR).
  \item[\keyterm{Frequency}]
    The rate at which loss events occur, typically expressed as number of
    events per unit of exposure per time period.
  \item[\keyterm{Key risk indicator (KRI)}]
    A metric that signals changes in risk exposure, enabling proactive
    monitoring rather than reactive response after losses occur.
  \item[\keyterm{Kurtosis}]
    A statistical measure of the fatness of distribution tails; high kurtosis
    indicates more frequent extreme events than the normal distribution
    predicts.
  \item[\keyterm{Loss distribution}]
    The probability distribution describing the likelihood of different loss
    amounts.
  \item[\keyterm{Mean}]
    The arithmetic average of a set of values.
  \item[\keyterm{Median}]
    The middle value when observations are arranged in order; half fall below
    and half above.
  \item[\keyterm{Monte Carlo simulation}]
    A computational technique using random sampling to simulate many possible
    outcomes and estimate probability distributions.
  \item[\keyterm{Normal distribution}]
    A symmetric, bell-shaped probability distribution characterized by mean
    and standard deviation; also called the Gaussian distribution.
  \item[\keyterm{Right-skewed distribution}]
    An asymmetric distribution with a long tail extending toward larger values;
    mean exceeds median.
  \item[\keyterm{Risk qualification}]
    Risk assessment using narrative descriptions, expert judgment, scenarios,
    and categorical ratings rather than numerical measurements.
  \item[\keyterm{Risk quantification}]
    Risk assessment using numerical data, statistical methods, and mathematical
    models to measure risk exposure.
  \item[\keyterm{Severity}]
    The magnitude of loss when an adverse event occurs, typically measured in
    dollars.
  \item[\keyterm{Skewness}]
    A statistical measure of the asymmetry of a probability distribution.
  \item[\keyterm{Standard deviation}]
    A measure of dispersion around the mean; the square root of variance,
    expressed in the same units as the data.
  \item[\keyterm{Stress testing}]
    Analysis of how risk exposures and outcomes would change under specified
    adverse scenarios.
  \item[\keyterm{Tail risk}]
    The risk of extreme losses occurring in the tails of a distribution,
    particularly the right tail for loss distributions.
  \item[\keyterm{Value at Risk (VaR)}]
    The maximum loss expected with a given confidence level over a specified
    time period; the threshold that losses will exceed with only small
    probability.
  \item[\keyterm{Variance}]
    A measure of dispersion around the mean, calculated as the average squared
    deviation from the mean.
\end{keyterms}

% -----------------------------------------------------------------------
\begin{shortanswerquestions}
\item A company experiences an average of 12 workplace accidents per year among
  300 employees, and each accident costs \$8,500 on average. Calculate:
  (a) frequency per employee per year; (b) expected annual loss; and
  (c) expected cost per employee per year.

\item An organization faces two risks: Risk~A has expected loss \$200,000 with
  standard deviation \$50,000; Risk~B has expected loss \$300,000 with standard
  deviation \$80,000. If the risks are independent (uncorrelated), what is the
  expected total loss? Estimate the standard deviation of total loss.

\item A company calculates that its VaR (95\%) for workers' compensation losses
  is \$450,000. Using the Northland Manufacturing case as a model, explain in
  plain English what this means and what it does \emph{not} tell you.

\item Referring to the Midtown Outfitters case (Section~\ref{sec:midtown-outfitters}),
  the company is deciding whether to increase its general liability insurance
  deductible from \$10,000 to \$25,000 per occurrence. Using the slip-and-fall
  data, what factors should inform this decision?
\end{shortanswerquestions}

% -----------------------------------------------------------------------
\begin{discussionquestions}
\item Explain why organizations cannot rely solely on managerial intuition for
  risk assessment. What problems does systematic risk quantification address?

\item Distinguish between risk quantification and risk qualification. When is
  each approach most appropriate? Provide a concrete example of a risk that
  would be best assessed primarily through qualitative methods.

\item A retailer's historical slip-and-fall claims show mean severity of
  \$12,000 and median severity of \$3,500. What does this tell you about the
  shape of the loss distribution? What does it imply for how the company should
  budget for, insure, and manage this risk?

\item In the Precision Parts case (Section~\ref{sec:precision-parts}), the risk
  manager noted positive correlation between equipment breakdowns and product
  defects. Explain why this correlation might exist and what it implies for risk
  management strategy.

\item Value at Risk is widely used but has significant limitations. Using the
  LTCM case as a reference, describe three important limitations of VaR and
  explain how risk managers should address each.

\item How does the Moneyball story illustrate the principles of risk
  quantification? What is the analogous challenge facing managers who rely on
  traditional qualitative judgment rather than systematic measurement?

\item Explain how quantitative risk measures (expected loss, VaR,
  frequency/severity statistics) integrate with qualitative risk assessment
  (scenarios, expert judgment, risk matrices) in a comprehensive ERM system.

\item A manufacturing company has never experienced a catastrophic fire at its
  main production facility but recognizes this as a significant potential risk.
  Why is quantitative analysis alone insufficient for assessing this risk? What
  qualitative approaches would complement limited quantitative data?

\item Choose a non-financial industry of your choice. Describe how the
  frequency-severity framework might be applied to a major risk that industry
  faces. What data would you collect? What decisions would the analysis inform?
  What limitations would you face?
\end{discussionquestions}

% -----------------------------------------------------------------------
\begin{minicase}{Workers' Compensation Risk Analysis: Northland Manufacturing}

Using the data described in Section~\ref{sec:northland-var} (Northland
Manufacturing), complete the following analysis in Excel.

\textbf{Part 1: Basic Calculations.} Create a spreadsheet with annual data for
Years~1--10 (columns for year, employees, work hours, injuries, and total
annual cost). Calculate total injuries and work hours over ten years, average
injuries per year, and average cost per year. Then compute frequency metrics:
injuries per 200,000 work hours; injuries per employee per year; and expected
injuries for a current workforce of 400 employees.

\textbf{Part 2: Severity Analysis.} Create a table showing the claim size
distribution (minor, moderate, and serious categories with counts, percentages,
and average costs). Calculate mean severity across all 138 claims and median
severity. Verify that the weighted average of your category means equals the
overall mean.

\textbf{Part 3: Expected Loss.} Calculate expected annual loss two ways:
(1) the mean of the historical annual cost column; and (2) frequency times
average severity. The two values should be very close---if they differ
substantially, check your frequency and severity calculations.

\textbf{Part 4: Value at Risk.} Sort annual losses from smallest to largest.
Calculate VaR (90\%) using
\excelcode{=PERCENTILE.INC(annual\_losses, 0.90)} and VaR (95\%) using
\excelcode{=PERCENTILE.INC(annual\_losses, 0.95)}. Identify which historical
year(s) exceeded each threshold. Create a histogram of annual losses with VaR
thresholds marked.

\textbf{Part 5: Interpretation.} In a text box or separate sheet, write
2--3 paragraphs addressing: What is the expected annual workers' compensation
cost for Northland? With 95\% confidence, what is the maximum annual cost the
company should budget for? If Northland wants to retain losses up to \$175,000
and insure above that amount, what percentage of historical years would have
triggered insurance coverage? What recommendations would you make to management
about insurance purchasing, safety programs, and budgeting?

\medskip
\textbf{Deliverable:} Submit your Excel workbook with clear labels, formulas
(not hard-coded values), and interpretation text.

\medskip
\textbf{Bonus Challenge (Optional):} Implement a simple Monte Carlo simulation.
Use \excelcode{RAND()} and a lookup table to simulate 1,000 years of losses.
For each simulated year, randomly determine the number of injuries (using the
historical rate) and randomly assign severities (using the empirical
distribution). Calculate VaR from your simulated distribution and compare to
VaR calculated from the ten-year historical data. Why might the two estimates
differ? Which do you trust more?

\end{minicase}

% -----------------------------------------------------------------------
\section*{References}
\addcontentsline{toc}{section}{References}
\begingroup
\sloppy

\begin{hangparas}{1.5em}{1}

Basel Committee on Banking Supervision. (2011). \emph{Basel~III: A global
regulatory framework for more resilient banks and banking systems.} Bank for
International Settlements. \url{https://www.bis.org/publ/bcbs189.pdf}

Centers for Disease Control and Prevention. (2023). \emph{National Healthcare
Safety Network (NHSN) overview.} U.S.\ Department of Health and Human
Services. \url{https://www.cdc.gov/nhsn/}

Centers for Medicare \& Medicaid Services. (2022). \emph{Hospital Readmissions
Reduction Program (HRRP).} U.S.\ Department of Health and Human Services.
\url{https://www.cms.gov/medicare/quality/value-based-programs/hrrp}

Lewis, M. (2003). \emph{Moneyball: The art of winning an unfair game.}
W.\,W.\ Norton \& Company.

Lowenstein, R. (2000). \emph{When genius failed: The rise and fall of
Long-Term Capital Management.} Random House.

\end{hangparas}

\endgroup
\end{document}
